\section{Optimization Problems using Function Maximization Oracle}\label{appendix:sec:Function-Maximization-Oracles}
In this section, we prove Lemmas 
\ref{lem:Preliminaries:wc-oracle-from-sc-oracle}, \ref{lemma:helper-for-p-max-oracle}, and \ref{lem:FunctionMaximizetionOracles:computing-maximal-pHw-tight-sets-with-p-max-oracle}. The proofs of these lemmas follow from known tools in the area of submodular functions. We present their proofs for the sake of completeness. 



\subsection{\functionMaximizationOracle{p} using \functionMaximizationOracleStrongCover{p}}\label{sec:wc-oracle-from-sc-oracle}
In this section, we prove \Cref{lem:Preliminaries:wc-oracle-from-sc-oracle}. 
\lemmaWCoraclefromSCoracle*
\begin{proof}
Let $u \in V$ be a vertex. We define the hypergraph $G^u := (V, E^u := \{e \cup \{u\} : e \in E\})$ and the weight function $w^u : E^u \rightarrow\Z_+$ as $w^u(e \cup \{u\}) := w(e)$ for every $e \in E$. We observe that $b_{(G, w)}(X) = d_{(G^u, w^u)}(X)$ for every set $X \subseteq V - \{u\}$. We also let $(Z^u, p(Z^u))$ denote the tuple obtained by querying \functionMaximizationOracleStrongCover{p}$((G_0, c_0) := (G^u, w^u), S_0 := S, T_0 := T\cup \{u\})$.

We now describe a procedure to compute an answer $(Z, p(Z))$ to the query $\functionMaximizationOracle{p}((G,w), S, T)$. We consider two cases to define our set $Z$. If $T \not= \emptyset$, then we choose an arbitrary vertex $t \in T$ and define $Z := Z^t$. Otherwise, $T = \emptyset$ and we define the set $Z$ to be an arbitrary optimum solution to $\max\{p(Z) - b_{(G, w)}(Z) : Z \in \calZ \}$, where $\calZ := \{Z^u: u \in V\} \cup \{ V\}$.  We return the tuple $(Z, p(Z))$ as the answer to the query $\functionMaximizationOracle{p}((G, w), S, T)$.

We now prove correctness of the above procedure. Let $Z$ be the set returned by the above procedure. Then, by definition $S\subseteq V\subseteq V-T$. Our goal is to show that $Z$ is an optimum solution to the following problem:
\begin{align}
    \max\left\{p(X)-b_{(G, w)}(X): S\subseteq X\subseteq V-T\right\} \label{eq:wc-oracle-from-sc-oracle}
\end{align}
Let $Y$ be an optimum solution to \eqref{eq:wc-oracle-from-sc-oracle}. It suffices to show that $p(Z)-b_{(G, w)}(Z) \ge p(Y)-b_{(G, w)}(Y) $.   
We consider two cases: suppose that $T\neq \emptyset$. Then, we have that 
\begin{align*}
p(Z)-b_{(G, w)}(Z)
&=p(Z^t)-d_{(G^t, w^t)}(Z^t)\\
&\ge p(Y)-d_{(G^t, w^t)}(Y) \\
&= p(Y) - b_{(G, w)}(Y) .
\end{align*} 
In the above, the first equality is because $Z=Z^t$, the inequality is because $Y\subseteq V-\{v\}$ and by definition of $Z^t$, and the last equality is because $b_{(G, w)}(Y) = d_{(G^v, w^v)}(Y)$ since $Y\subseteq V-\{t\}$. 
Next, suppose that $T=\emptyset$. We consider two subcases. Firstly, suppose that $V$ is the unique optimum solution to \eqref{eq:wc-oracle-from-sc-oracle}. Then, $Z=Y=V$ by definition and we are done. 
Secondly, suppose that $Y\neq V$. Let $v\in V-Y$.  
Then, we have the following:
\begin{align*}
p(Z)-b_{(G, w)}(Z)
%&=p(Z^u)-b_{(G, w)}(Z^u)\\
&\ge p(Z^v) - b_{(G, w)}(Z^v) \\
&=p(Z^v)-d_{(G^v, w^v)}(Z^v) \\
&\ge p(Y)-d_{(G^v, w^v)}(Y) \\
&=p(Y)-b_{(G, w)}(Y).
\end{align*} 
In the above sequence, the first inequality is because of the choice of the set $Z$ (i.e., $Z$ is an optimum solution to $\max\{p(Z)-b_{(G, w)}(Z): Z\in \calZ\}$ and $Z^v \in \calZ$), the first equality is because  $b_{(G, w)}(Z^v) = d_{(G^v, w^v)}(Z^v)$ since $Z^v\subseteq V-\{v\}$, the second inequality is by definition of $Z^v$, and the second equality is because $b_{(G, w)}(Y) = d_{(G^v, w^v)}(Y)$ since $Y\subseteq V-\{v\}$. Thus, all inequalities in the above sequence are equalities, and we have that $p(Z)-b_{(G, w)}(Z) = p(Y)-b_{(G, w)}(Y)$. Thus, in both cases, the set $Z$ returned by our procedure is an optimum solution to \eqref{eq:wc-oracle-from-sc-oracle}.

We now analyze the runtime of the procedure and the number of \functionMaximizationOracleStrongCover{p} queries. Let $u \in V$ be a vertex. Constructing the hypergraph $(G^{u}, w^{u})$ takes $O(|V| + |E|)$. Thus, a tuple $(Z^{u}, p(Z^u))$ takes $O(|V|+|E|)$ time and one query to \functionMaximizationOracleStrongCover{p}. The runtime includes the time to construct the inputs to the query. Moreover, the input hypergraph to the query has $|V|$ vertices and at most $|E|$ hyperedges. Thus, if $T \not = \emptyset$, the procedure makes a single \functionMaximizationOracleStrongCover{p} query and runs in $O(|V| + |E|)$ time. Alternatively, if  $T  = \emptyset$, then the procedure makes $|V| + 1$ \functionMaximizationOracleStrongCover{p} queries and runs in $O(|V|(|V| + |E|))$ time. In both cases, the runtime includes the time to construct the inputs for the \functionMaximizationOracleStrongCover{p} queries. Moreover, the input hypergraph to each query has $|V|$ vertices and at most $|E|$ hyperedges.
\end{proof}

\subsection{Skew-supermodularity of a function \(p\)}\label{sec:helper-for-p-max-oracle}
In this section, we prove \Cref{lemma:helper-for-p-max-oracle}. 
\lempmaxoracle*
\begin{proof}
    For ease of notation, let $p=p_{(G, c, r)}$. 
    The function $R:2^V\rightarrow \Z_{\ge 0}$ is symmetric skew-supermodular (see Lemma 8.1.9 of \cite{Frank-book}) and the function $d_{(G,c)}:2^V\rightarrow \Z_{\ge 0}$ is symmetric submodular. The function $p$ is obtained by subtracting a symmetric submodular function from a symmetric skew-supermodular and is hence, symmetric skew-supermodular. 

    Next, we show that \functionMaximizationOracleStrongCover{p} can be implemented in strongly polynomial-time. 
    Let the input to the oracle be the hypergraph $(G_0=(V, E_0), c_0: E_0\rightarrow \Z_+)$ and disjoint subsets $S_0, T_0\subseteq V$. 
    We construct the hypergraph $(G+G_0, c+c_0)$. Moreover, for distinct vertices $u, v\in V$ with $u\in V-T_0$ and $v\in V-S_0$, we compute  
    \begin{align*}
    \lambda_{(G+G_0, c+c_0)}(u, v)&\coloneqq \min\left\{d_{(G+G_0, c+c_0)}(Z): S_0\cup \{u\}\subseteq Z\subseteq V-(T_0\cup\{v\})\right\} 
    \end{align*}
    and the set $Z(u, v)$ that achieves the minimum for $\lambda_{(G+G_0, c+c_0)}(u, v)$. We note that $\lambda_{(G+G_0, c+c_0)}(u, v)$ and $Z(u,v)$ can be computed using a single call to a min $(s,t)$-cut oracle on the hypergraph $(G+G_0, c+c_0)$. This hypergraph has $|V|$ vertices and at most $|E_0|+|E|$ hyperedges.  

    To implement the \functionMaximizationOracleStrongCover{p}, we observe that 
    \begin{align*}
    \max&\left\{p(Z)-d_{(G_0, c_0)}(Z): S_0\subseteq Z\subseteq V-T_0\right\}\\
    &=\max\left\{R(Z)-d_{(G, c)}(Z)-d_{(G_0, c_0)}(Z): S_0\subseteq Z\subseteq V-T_0\right\}\\
    &=\max\left\{r(u, v)-\lambda_{((G, c), (G_0, c_0))}(u, v): u\in V-T_0, v\in V-S_0, u\neq v\right\}. 
    \end{align*}
    The RHS problem can be solved by iterating over all distinct vertices $u\in V-T_0$ and $v\in V-S_0$ and returning a pair for which the objective function $r(u, v)-\lambda_{((G, c), (G_0, c_0))}(u, v)$ is maximum. Let $u\in V-T_0$ and $v\in V-S_0$ be a pair of vertices that achieves the maximum in the RHS problem. Then, we observe that the set $Z(u, v)$ is a maximizer for the LHS problem. Thus, it suffices to return $Z(u, v)$ and the value $p(Z(u, v))$. We observe that the value $p(Z(u, v))$ can be computed by subtracting $d_{(G_0, c_0)}(Z(u,v))$ from the objective value of the RHS problem.

    The above procedure can be implemented in $O(|V|^2(|V| + |E| + |E_0|))$ time and $O(|V|^2)$ calls to a min $(s,t)$-cut oracle on hypergraphs with $|V|$ vertices and $|E| + |E_0|$ hyperedges. Then, the claimed runtime is obtained by using a crude upper bound on the run-time of standard deterministic min $(s,t)$-cut algorithms for hypergraphs---we recall that min $(s,t)$-cut in a $n$-vertex $m$-hyperedge hypergraph can be solved in $O((m+n)mn)$.
\end{proof}


\subsection{Computing the Family of Maximal \((p, H, w)\)-Tight Sets}\label{appendix:sec:Function-Maximization-Oracles:computing-maximal-pHw-tight-sets-family}
In this section, we prove \Cref{lem:FunctionMaximizetionOracles:computing-maximal-pHw-tight-sets-with-p-max-oracle}. 
We recall that for a set function $p:2^V\rightarrow\Z$ and hypergraph $(H = (V, E), w)$, a set $X\subseteq V$ is $(p, H, w)$-tight if $p(X) = b_{(H,w)}(X)$. The family $\maximalTightSetFamily{p, H, w}$ denotes the family of all inclusion-wise maximal $(p, H, w)$-tight sets. In this section, we show that the family $\maximalTightSetFamily{p, H, w}$ can be computed in $O(|V|^2(|V| + |E|))$ time using $O(|V|^2)$ queries to \functionMaximizationEmptyOracle{p} if the function $p$ and hypergraph $(H, w)$ satisfy certain properties.

\lemComputingMaximalTightSetFamily*
\begin{proof}
We describe a procedure that recursively constructs \maximalTightSetFamily{p, H, w}. The procedure finds a $(p, H, w)$-tight set $X\subseteq V$, and then searches for a maximal $(p, H, w)$-tight set $Y$ containing the set $X$ --- we describe this computation in the paragraph below. 
The procedure is recursive and it reaches the base case if 
the computed set $X$ is not a $(p, H, w)$-tight set or $X,Y = \emptyset$. If $X$ is not a $(p, H, w)$-tight set, then the procedure returns the empty set, and if $X,Y = \emptyset$, then the procedure returns $\{\emptyset\}$.
If both of the base case conditions are not satisfied, the procedure recurses on the function $p':2^{V'}\rightarrow\Z$ defined on the ground set $V':=V - Y$ as $p'(Z) = p(Z)$ for every $Z \subseteq V'$, and the hypergraph $(H' = (V', E'), w')$ where $E' := \{e \cap V' : e \in E\}$ and $w'(e'):= \sum_{e \in E}w(e)\cdot\indicator_{e\cap V' = e'}$ for every $e' \in E'$, to obtain the family $\maximalTightSetFamily{p', H', w'}$.
It then returns the family $\{Y\} \cup (\maximalTightSetFamily{p', H', w'} - \{\emptyset\})$. 
The procedure terminates within $|V|$ recursive calls because the size of the ground set decreases during every recursive call (i.e. $|V'| < |V|$). Furthermore, the correctness of the procedure follows from disjointness of the family $\maximalTightSetFamily{p, H, w}$ by \Cref{lem:weak-to-strong:uncrossing:maximal-tight-set-family-disjoint} and because $b_{(H,w)}(X) = b_{(H', w')}(X)$ for each $X \subseteq V'$, where $(H', w')$ is the hypergraph constructed as input to the subsequent recursive call.
    
We now bound the run-time and the number of queries required by the procedure described above. Consider a recursive call of the procedure.
We note that a $(p, H, w)$-tight set exists if and only if the set returned by the query $\functionMaximizationOracle{p}((G_0, c_0):= (H, w), S_0:=\emptyset, T_0:=\emptyset)$ is a $(p, H, w)$-tight set.
Furthermore, verifying whether the returned set $X$ is $(p, H, w)$-tight can be done in $O(|V| + |E|)$ time by checking whether $p(X) = b_{(H,w)}(X)$.
Thus, a $(p, H, w)$-tight set $X\subseteq V$ can be found in  $O(|V| + |E|)$ time using one query to \functionMaximizationOracle{p}. Moreover, the input hypergraph to \functionMaximizationOracle{p} query has $|V|$ vertices and $|E|$ hyperedges.
Next, \Cref{claim:FunctionMaximizetionOracles:maximal-tight-set-containing-given-tight-set} below shows that we can find a maximal $(p, H, w)$-tight set $Y$ containing the set $X$ in $O(|V|(|V|+|E|))$ time using $O(|V|)$ queries to \functionMaximizationOracle{p} --- here, the runtime includes the time required to construct the inputs to \functionMaximizationOracle{p}; moreover, each query to \functionMaximizationOracle{p} takes the hypergraph $(H, w)$ as input.
Let $p':2^{V'}\rightarrow\Z$ be the input function to the subsequent recursive call. Let $(G = (V', E_G), c)$ be a hypergraph such that $|E_G| \leq |E|$ and $S, T \subseteq V'$ be disjoints sets. 
Then, \Cref{claim:FunctionMaximizetionOracles:computing-maximal-tight-sets:maximization-oracle-for-restriction} below shows that a query to $\functionMaximizationOracle{p'}((G, c), S, T)$ can be implemented in $O(|V| + |E|)$ time using one query to \functionMaximizationOracle{p} --- here, the runtime includes the time required to construct the input to \functionMaximizationOracle{p} query; moreover the hypergraph input to \functionMaximizationOracle{p} query has $|V|$ vertices and $|E_G|$ hyperedges. We note that all hypergraphs $(G = (V', E_G), c)$ that are used as inputs to \functionMaximizationOracle{p'} during subsequent recursive calls satisfy $|E_G| \leq |E|$.
Thus, every recursive call of the procedure can be implemented to run in $O(|V|(|V|+|E|))$ time and make $O(|V|)$ calls to \functionMaximizationOracle{p}, where the runtime includes the time to construct the inputs to each query to \functionMaximizationOracle{p}. Consequently, the procedure runs in time $O(|V|^2(|V|+|E|))$ and makes $O(|V|^2)$ queries to $\functionMaximizationOracle{p}$, where the runtime includes the time to construct the inputs to the \functionMaximizationOracle{p} queries.
\end{proof}

\begin{claim}\label{claim:FunctionMaximizetionOracles:maximal-tight-set-containing-given-tight-set}
         Let $X \subseteq V$ be a $(p, H, w)$-tight set. Then, a maximal $(p, H, w)$-tight set containing the set $X$ can be computed in $O(|V|(|V|+|E|))$ time using $O(|V|)$ queries to \functionMaximizationOracle{p}. The runtime includes the time required to construct the inputs to \functionMaximizationOracle{p}. Moreover, each query to \functionMaximizationOracle{p} takes the hypergraph $(H, w)$ as input.
\end{claim}
\begin{proof}
We describe a recursive procedure to compute a maximal $(p, H, w)$-tight set containing the set $X$. For every $u \in V - X$, let $Z^u$ be the set returned by querying $\functionMaximizationOracle{p}((H_0, w_0):= (H, w), S_0 := X\cup\{u\}, T_0 := \emptyset)$. If there exists a vertex $u \in X$ such that the set $Z^u$ is a $(p, H, w)$-tight set, then recursively compute a $(p, H, w)$-tight set containing the set $X':= Z^u$. Otherwise, return the set $X$. 

The correctness of the above procedure is because of the following. The above procedure terminates (in at most $|V|$ recursive calls) because the cardinality of the input set $X$ increases during each recursive call (i.e., $|X'| > |X|$). Let $Y \subseteq V$ be a maximal $(p, H, w)$-tight set such that $X \subseteq Y$. If $Y = X$, then there are no $(p, H, w)$-tight sets that strictly contain the set $X$ and the procedure returns the set $X$. Suppose that $X \subsetneq Y$ and let $u \in Y - X$. Then, there exists a $(p, H, w)$-tight set $Z$ such that $X\cup\{u\} \subsetneq Z \subseteq Y$. Thus, the set $Z^u$ computed by the procedure will be a $(p, H, w)$-tight set that strictly contains the set $X$. Using this fact, the procedure can be shown to return a maximal $(p, H, w)$-tight set containing the set $X$ by induction on $|X|$.   
    
We now describe a slight modification of the above procedure that leads to an improved runtime and a smaller number of queries to \functionMaximizationOracle{p}. Let $u \in V - X$ be a vertex. If the set $Z^{u}$ is $(p, H, w)$-tight, then we include $u$ in the input set $S_0$ to all $\functionMaximizationOracle{p}$ queries in subsequent recursive calls of the procedure. Alternatively, if the set $Z^{u}$ is not $(p, H, w)$-tight, then we include $u$ in the input set $T_0$ to all $\functionMaximizationOracle{p}$ queries in subsequent recursive calls procedure. Consequently, the number of queries to $\functionMaximizationOracle{p}$ (and number of recursive calls of the procedure) is at most $|V|$. We note that constructing the input to each query and verifying whether the set returned by the query is $(p, H, w)$-tight can be done in $O(|V| + |E|)$ time. Thus, the procedure terminates in $O(|V|(|V|+|E|))$ time. Here, the runtime includes the time to construct the inputs to the \functionMaximizationOracle{p} queries. Moreover, each query takes the hypergraph $(H,w)$ as input. 
\end{proof}
    
\begin{claim}\label{claim:FunctionMaximizetionOracles:computing-maximal-tight-sets:maximization-oracle-for-restriction}
    Let $p:2^V\rightarrow\Z$ be a function. Let $Y \subseteq V$ and let $p'$ denote the function $p$ restricted to the ground set $V':=V - Y$. Then, for a given hypergraph $(G = (V', E_G), c:E_G\rightarrow\Z_+)$ such that $|E_G| \leq |E|$ and disjoint subsets $S, T \subseteq V'$, $\functionMaximizationOracle{p'}((G,c), S, T, y)$ can be computed in $O(|V| + |E|)$ time using one query to $\functionMaximizationOracle{p}$.
    The runtime includes the time required to construct the input to \functionMaximizationOracle{p} query. Moreover, the hypergraph input to \functionMaximizationOracle{p} query has $|V|$ vertices and $|E_G|$ hyperedges.
\end{claim}
\begin{proof}
    Let $(G_1, c_1)$ be the hypergraph obtained from $(G,c)$ by adding the vertices $Y$. Then, the answer to the query $\functionMaximizationOracle{p'}((G, c), S, T)$ is the same as the answer to the query $\functionMaximizationOracle{p}((G_1, c_1), S, T\cup Y)$.
\end{proof}





